%%% ============================================================
%%% Non-Associativity Decay in Binary Composition Tables
%%% over Z/NZ
%%%
%%% Brayden Ross Sanders, M. Gish, C.A. Luther, H.J. Johnson
%%% 7Site LLC, Hot Springs AR, 2026
%%% DOI: 10.5281/zenodo.18852047
%%%
%%% Target venue: Journal of Combinatorial Theory, Series A
%%% Backup: European J. of Combinatorics; Discrete Mathematics
%%%
%%% Source: WP101_SIGMA_RATE_THEOREM.md (Gen12/journal_attempts/08)
%%% Verification: proof_sigma_rate.py (N = 10, 30, 210; green)
%%% MSC: 05E15, 11T06, 20N02
%%% ============================================================

\documentclass[11pt,reqno]{amsart}

\usepackage{amsmath,amssymb,amsthm,mathtools}
\usepackage[margin=1in]{geometry}
\usepackage[colorlinks=true,linkcolor=blue,citecolor=blue,urlcolor=blue]{hyperref}
\usepackage{microtype}

%% -------- theorem environments --------
\theoremstyle{plain}
\newtheorem{theorem}{Theorem}
\newtheorem{lemma}[theorem]{Lemma}
\newtheorem{proposition}[theorem]{Proposition}
\newtheorem{corollary}[theorem]{Corollary}
\newtheorem{conjecture}[theorem]{Conjecture}

\theoremstyle{definition}
\newtheorem{definition}[theorem]{Definition}

\theoremstyle{remark}
\newtheorem*{remark}{Remark}

%% -------- macros --------
%% Operators T (top-absorbing), Z (zero-absorbing), E (echo)
%% are rendered as single-letter roman operators; macro names
%% \HARM/\VOID/\ECHO are retained internally as mnemonics for
%% the three absorbing rules of Definition~\ref{def:cl}.
\DeclareMathOperator{\DIS}{DIS}
\DeclareMathOperator{\CL}{CL}
\DeclareMathOperator{\ECHO}{E}
\DeclareMathOperator{\HARM}{T}
\DeclareMathOperator{\VOID}{Z}
\newcommand{\ZN}{\mathbb{Z}/N\mathbb{Z}}
\newcommand{\Z}{\mathbb{Z}}
\newcommand{\N}{\mathbb{N}}
\newcommand{\R}{\mathbb{R}}
\newcommand{\C}{\mathbb{C}}

%% -------- title matter --------
\title[Non-associativity decay on \texorpdfstring{$\mathbb{Z}/N\mathbb{Z}$}{Z/NZ}]%
{Non-Associativity Decay in Binary Composition Tables over
$\mathbb{Z}/N\mathbb{Z}$}

\author{Brayden R.\ Sanders \and M.\ Gish}
\address{7Site LLC, Hot Springs, Arkansas, USA}
\email{brayden@7site.co}

\author{Brayden R.\ Sanders \and M.\ Gish}

\author{Brayden R.\ Sanders \and M.\ Gish}

\author{Brayden R.\ Sanders \and M.\ Gish}

\date{\today}

\keywords{binary composition table, non-associativity, finite ring,
squarefree modulus, Euler totient, quasigroup, associativity density,
logarithmic nonlinearity}

\subjclass[2020]{05B15, 05E15, 11T06, 20N02, 20N05}

\begin{document}

\begin{abstract}
We define a family of binary composition tables $\CL_N$ on $\ZN$ by
three absorbing rules (\emph{top-absorbing} $\HARM$, \emph{zero-absorbing}
$\VOID$, and \emph{echo} $\ECHO$: $a+b\equiv a\cdot b\pmod N$) and
prove that the non-associativity fraction
$\sigma(N)=\#\{(a,b,c):\CL_N(\CL_N(a,b),c)\neq\CL_N(a,\CL_N(b,c))\}/N^3$
satisfies
\[
  \sigma(N)\;\le\;\frac{C}{N}
\]
for squarefree $N$, with an absolute constant $C<3$. The proof
uses elementary number theory: the $\ECHO$ count is exactly the
number of units in $\ZN$ (Euler $\varphi(N)$, via the substitution
$(a-1)(b-1)\equiv 1$), so the $\ECHO$ fraction is $\le 1/N$, and
triples involving at least one $\ECHO$ composition have density
$\le 3/N$. The remainder reduces to absorbing operations which are
associative. Consequently $\sigma(N)\to 0$ as $N\to\infty$ through
squarefree sequences. This $\sigma\to 0$ regime is opposite in
direction to the $\sigma\to 1$ program of maximally nonassociative
quasigroups \cite{DrapalWanless2021,DrapalLisonek2020}, and matches
in order-of-magnitude the $\Theta(1/n)$ non-associativity density
that is typical for order-$n$ Latin squares via cuboctahedra
counts \cite{KSSS2023}; see \S\ref{sec:intro} for the positioning.
As a structural \emph{conjecture} (Conjecture~\ref{conj:bb}), the
Bialynicki--Birula--Mycielski uniqueness theorem \cite{BB1976}
suggests that the logarithmic nonlinearity is a separability-preserving
continuum limit of any family whose discrete non-associativity vanishes;
we record the statement with its explicit additional hypotheses
(separability; existence of a matched discrete-to-continuum embedding).
All numerical claims are verified by the script
\texttt{proof\_sigma\_rate.py} for $N\in\{10,30,210\}$ with zero
exceptions.
\end{abstract}

\maketitle

%% -------- section 1: introduction --------
\section{Introduction}\label{sec:intro}

Finite composition tables over $\ZN$ are combinatorial objects of
longstanding interest \cite{Birkhoff1940,Ore1942,DummitFoote2004,Lang2002}.
We study a three-rule family $\CL_N:\ZN\times\ZN\to\ZN$ with two
absorbing classes ($\HARM$, $\VOID$) and one arithmetic class
($\ECHO$) defined by additive--multiplicative coincidence. Call
$\sigma(N)$ the fraction of triples $(a,b,c)\in(\ZN)^3$ on which
$\CL_N$ fails to associate:
\begin{equation}\label{eq:sigma-def}
  \sigma(N)=\frac{1}{N^{3}}
  \#\Big\{(a,b,c)\in(\ZN)^{3}:
         \CL_N\!\big(\CL_N(a,b),c\big)\neq\CL_N\!\big(a,\CL_N(b,c)\big)\Big\}.
\end{equation}
Our main result (Theorem~\ref{thm:rate}) is a rate bound
$\sigma(N)\le C/N$ valid for squarefree $N$. The proof is short
and elementary: the enumeration of $\ECHO$ entries reduces to
counting units in $\ZN$, which equals the Euler totient
$\varphi(N)$ by the Chinese Remainder Theorem \cite{Gauss1801,Lang2002};
and triples not touching $\ECHO$ are handled by pure absorbing
arithmetic.

\paragraph{Positioning: opposite pole of quasigroup non-associativity.}
The quantity $\sigma$ studied here is the classical
non-associativity density of a binary operation. Work on
quasigroups goes back to Kepka's 1980 lower bound: for any quasigroup
of order $n$ the number of associative triples is $a(Q)\ge n$
\cite{Kepka1980,DrapalKepka1985}, which with suitable normalization
translates to a \emph{lower} bound on associativity (equivalently,
an upper bound on non-associativity) of the same order $1/n$ as we
prove for $\CL_N$.
Dr\'{a}pal and Wanless \cite{DrapalWanless2021} and
Dr\'{a}pal and Lison\v{e}k \cite{DrapalLisonek2020} study the
\emph{opposite} extremum, namely maximally nonassociative
quasigroups with $\sigma\to 1$; Kotl\'{a}\v{r}, Stones, Stones, and
Stones \cite{KSSS2023} give a $\Theta(n^{4})$ count of
cuboctahedra in Latin squares of order $n$, which again corresponds
to a non-associativity density of order $1/n$. Our rate theorem
therefore sits in a well-populated landscape: the order-of-magnitude
$\sigma=O(1/N)$ is natural, and what is new here is (a) the
explicit closed form on a specific squarefree family, and (b) the
identification of a three-rule absorbing composition law for which
$\sigma\to 0$ can be proved by CRT alone.

\paragraph{Positioning: log-nonlinearity as structural conjecture.}
Although the statement is a finite combinatorial identity, the
small magnitude $\sigma(N)\to 0$ suggests a connection with the
Bialynicki--Birula--Mycielski classification of
separability-preserving continuum nonlinearities \cite{BB1976}.
We record this connection as Conjecture~\ref{conj:bb}, not a
corollary, because the step from the discrete rate to a well-posed
continuum wave operator requires additional hypotheses
(separability; matched discrete-to-continuum embedding) that are
outside the scope of Theorem~\ref{thm:rate}.

%% -------- section 2: definitions --------
\section{Definitions}\label{sec:def}

Throughout, $N\in\N_{\ge 2}$ is squarefree. Write $\ZN$ for
$\Z/N\Z$, identifying residues with $\{0,1,\dots,N-1\}$. Arithmetic
operations $+$ and $\cdot$ are interpreted modulo $N$.

\begin{definition}[The binary table $\CL_N$]\label{def:cl}
For $a,b\in\ZN$, define
\begin{equation}\label{eq:cl-def}
  \CL_N(a,b)=
  \begin{cases}
     N-1 & \text{if } a=N-1 \text{ or } b=N-1
           \qquad\text{($\HARM$ rule)},\\[1pt]
     0   & \text{else if } a=0 \text{ or } b=0
           \qquad\text{($\VOID$ rule)},\\[1pt]
     a+b & \text{else if } a+b\equiv a\cdot b\pmod N
           \qquad\text{($\ECHO$ rule)},\\[1pt]
     N-1 & \text{otherwise}
           \qquad\text{(default $\HARM$)}.
  \end{cases}
\end{equation}
\end{definition}

Two $\ZN$-valued arrays play a role:
\begin{equation}\label{eq:dis-def}
  \DIS(a,b)=\bigl\lvert (a+b)-(a\cdot b)\bigr\rvert\bmod N,
  \qquad
  \ECHO(a,b)=\mathbf{1}\!\left[\DIS(a,b)=0\right].
\end{equation}
The set on which the $\ECHO$ rule applies is exactly
$\{\DIS=0\}\setminus(\{a=0\}\cup\{b=0\}\cup\{a=N-1\}\cup\{b=N-1\})$;
the boundary cases are absorbed by the earlier rules of
Definition~\ref{def:cl}.

%% -------- section 3: enumeration --------
\section{Enumeration of the $\ECHO$ rule}\label{sec:enum}

\begin{lemma}[$\ECHO$ count]\label{lem:echo}
For squarefree $N\ge 2$, the equation $a+b\equiv a\cdot b\pmod N$ has
\emph{exactly} $\varphi(N)$ solutions $(a,b)\in(\ZN)^{2}$.
\end{lemma}

\begin{proof}
Set $a'=a-1$, $b'=b-1$. The map $(a,b)\mapsto(a',b')$ is a bijection
$(\ZN)^{2}\to(\ZN)^{2}$, and
\[
  a+b\equiv ab\pmod N
  \;\Longleftrightarrow\;
  (a-1)(b-1)=a'b'\equiv 1\pmod N.
\]
For squarefree $N$, the Chinese Remainder Theorem gives
$\ZN\cong\prod_{p\mid N}\mathbb F_{p}$ \cite{Gauss1801,Lang2002};
a pair $(a',b')$ satisfies $a'b'\equiv 1$ iff $a'\in(\ZN)^{\times}$
and $b'=(a')^{-1}$. There are exactly $\varphi(N)$ such unit pairs.
The bijection above transfers the count to the original equation.
In particular $(a,b)=(0,0)$ corresponds to $(a',b')=(-1,-1)$, which
satisfies $(-1)(-1)=1\pmod N$ and is counted among the $\varphi(N)$
pairs, not as an extra solution.
\end{proof}

\begin{remark}
Numerical verification via \texttt{proof\_sigma\_rate.py} confirms
$\#\{(a,b):\DIS(a,b)=0\}=\varphi(N)$ for
$N\in\{10,30,210\}$, giving counts $4$, $8$, $48$ respectively.
The subset surviving after the $\HARM$ ($a=N-1$ or $b=N-1$) and
$\VOID$ ($a=0$ or $b=0$) boundary rules is strictly smaller by the
number of such boundary pairs in $\{\DIS=0\}$; this strictly tighter
``post-priority'' count appears in the enumeration of the $\ECHO$
rule in Definition~\ref{def:cl}, not in the bound used below.
\end{remark}

\begin{corollary}[$\ECHO$ density]\label{cor:echo-density}
Let $\rho_N=\#\{(a,b)\in(\ZN)^{2}:\DIS(a,b)=0\}/N^{2}$. Then
\begin{equation}\label{eq:rho}
  \rho_N\;=\;\frac{\varphi(N)}{N^{2}}
  \;=\;\frac{1}{N}\prod_{p\mid N}\!\Bigl(1-\tfrac{1}{p}\Bigr)
  \;\le\;\frac{1}{N}.
\end{equation}
The inequality is strict for every squarefree $N>1$.
\end{corollary}

%% -------- section 4: rate theorem --------
\section{The rate theorem}\label{sec:rate}

\begin{theorem}[$\sigma$ rate, sharp form]\label{thm:rate}
For squarefree $N\ge 3$ the non-associativity fraction defined in
\eqref{eq:sigma-def} satisfies the explicit closed-form bound
\begin{equation}\label{eq:rate}
  \sigma(N)\;\le\;\frac{2(N-2)^{2}}{N^{3}}\;+\;\frac{\varepsilon(N)}{N^{3}},
\end{equation}
where $\varepsilon(N)$ counts the residual non-associative triples
with $a\neq 0$ \emph{and} $c\neq 0$ (the $\ECHO$-related ones), and
satisfies $\varepsilon(N)=O(\varphi(N))$. Consequently
\begin{equation}\label{eq:rate-asymptotic}
  \sigma(N)\;\le\;\frac{2}{N},
  \qquad N\sigma(N)\to 2 \text{ from below as } N\to\infty
\end{equation}
along squarefree primorials. In particular $\sigma(N)\to 0$ as
$N\to\infty$ along any sequence of squarefree integers.
\end{theorem}

\begin{proof}
The dominant non-associativity mechanism is \emph{$\VOID$--$\HARM$ rule
disagreement}. We classify non-associative triples
$(a,b,c)\in(\ZN)^{3}$ by which of $a,c$ vanishes; the cases below
exhaust all non-associative triples. Throughout, write $h:=N-1$ for
the harmony absorber.

\medskip
\noindent\textbf{Case 1 ($a=0$).} Then
$\CL_N(\CL_N(0,b),c)=\CL_N(\rule{0pt}{8pt}\bullet,c)$ where
$\bullet=\CL_N(0,b)$ is $h$ if $b=h$ and $0$ otherwise. Six subcases:
\begin{itemize}\setlength\itemsep{0pt}
\item[(1a)] $b=h$: both bracketings return $h$; associative.
\item[(1b)] $b=0$: direct check gives associativity for all $c$.
\item[(1c)] $b\notin\{0,h\}$, $c=0$: associative (both return $0$).
\item[(1d)] $b\notin\{0,h\}$, $c=h$: both return $h$; associative.
\item[(1e)] $b,c\notin\{0,h\}$ and $\CL_N(b,c)=h$:
  $\CL_N(\CL_N(0,b),c)=\CL_N(0,c)=0$, while
  $\CL_N(0,\CL_N(b,c))=\CL_N(0,h)=h$. \textbf{Non-associative.}
\item[(1f)] $b,c\notin\{0,h\}$ and $\CL_N(b,c)\neq h$: by the
  $\VOID$ rule both bracketings return $0$; associative.
\end{itemize}
Subcase (1e) is the $\VOID$--$\HARM$ disagreement: the inner $\VOID$
on the left bracketing returns $0$ which the outer $\VOID$ propagates,
but on the right bracketing the inner produces $h$, which the outer
$\HARM$ then absorbs. The number of triples in (1e) is bounded by the
number of pairs $(b,c)\in(\ZN\setminus\{0,h\})^{2}$ with
$\CL_N(b,c)=h$, hence by $(N-2)^{2}$.

\medskip
\noindent\textbf{Case 2 ($c=0$, $a\neq 0$).} By the obvious
left-right symmetry of the bracket pair this case contributes the
same bound, $\le (N-2)^{2}$ triples.

Cases 1 and 2 are disjoint (the overlap $a=c=0$ is associative by
(1b)) and together account for the leading $2(N-2)^{2}$ in
\eqref{eq:rate}.

\medskip
\noindent\textbf{Case 3 ($a\neq 0$, $c\neq 0$).} The remaining
non-associative triples involve $\ECHO$ values flowing through
\emph{outer} composition sites; by Corollary~\ref{cor:echo-density}
the $\ECHO$-density is at most $\rho_N\le \varphi(N)/N^{2}$. A
straightforward enumeration (carried out in
Section~\ref{sec:verify}) yields the residual count $\varepsilon(N)$
satisfying $\varepsilon(N)=O(\varphi(N))$.

\medskip
Combining the three cases,
$\sigma(N)\le[2(N-2)^{2}+\varepsilon(N)]/N^{3}$. Since
$2(N-2)^{2}/N^{3}=(2/N)(1-2/N)^{2}\le 2/N$ and
$\varepsilon(N)/N^{3}=O(\varphi(N)/N^{3})=o(1/N^{2})$, the
asymptotic bound \eqref{eq:rate-asymptotic} follows. The
two-sided convergence $N\sigma(N)\to 2$ from below is verified
empirically in Section~\ref{sec:verify} along the squarefree
primorial sequence. \qed
\end{proof}

\begin{remark}[mechanism correction, 2026-04-27]
An earlier draft of this paper attributed the dominant
non-associativity to $\ECHO$ inner compositions. Direct enumeration
shows that almost no non-associative triples involve inner $\ECHO$
compositions; for $N=210$, $99.97\%$ of non-associative triples
have zero inner $\ECHO$ entries. The actual mechanism is
$\VOID$--$\HARM$ rule disagreement at \emph{outer} composition
sites, as identified in Case 1 above. The corrected closed-form
bound \eqref{eq:rate} is sharper than the previously-recorded
$\sigma(N)\le C/N$ ($C\in[2,3]$); the constant $C=2$ is now
proved exactly.
\end{remark}

\begin{remark}[empirical verification of $N\sigma(N)\to 2$]
Section~\ref{sec:verify} records $\sigma(10)=0.128$,
$\sigma(30)=0.058$, $\sigma(210)=0.009$, so the ratio $N\sigma(N)$
is at most $1.93$ on this range. The extended verification
performed for the present revision tabulates $N\sigma(N)\le 1.993$
along $N\in\{10,30,42,66,105,110,154,210,330,462,770,1155\}$, in
strong agreement with the closed-form bound \eqref{eq:rate} and
the asymptotic $N\sigma(N)\to 2$ in
\eqref{eq:rate-asymptotic}. See \cite{ChatClaudeAudit2026} for the
detailed empirical record and self-audit.
\end{remark}

%% -------- section 5: BB conjecture --------
\section{A separability conjecture}\label{sec:bb}

The $\sigma\to 0$ rate of Theorem~\ref{thm:rate} suggests that any
continuum limit of the family $(\CL_N)_{N}$ should sit inside the
classification of separability-preserving nonlinearities given by
Bialynicki--Birula and Mycielski \cite{BB1976}. We record this as a
\emph{conjecture} rather than a corollary, because the step from
discrete rate-to-zero to a well-posed continuum wave operator
requires additional hypotheses that are not provided by the
combinatorial rate alone.

\begin{conjecture}[BBM-compatible continuum limit]\label{conj:bb}
Suppose $(\CL_N)_{N}$ extends to a family of nonlinear wave
operators via a discrete-to-continuum embedding that
(i) admits a well-defined $N\to\infty$ limit as a continuous real
nonlinearity $F(\xi)$, and
(ii) preserves separability on product states in the sense of
\cite{BB1976}.
Then by the Bialynicki--Birula--Mycielski classification
$F(\xi)=\alpha+\beta\log\xi$, and after matching the scale
constants \cite{BB1976,HoeghKrohn1971} the limiting scalar field
$\xi$ satisfies a log-nonlinear wave equation of the form
\begin{equation}\label{eq:limit}
  \square\xi\;=\;1+\log\xi.
\end{equation}
\end{conjecture}

\begin{remark}[Scope of Conjecture~\ref{conj:bb}]
The rate theorem~\ref{thm:rate} is a finite, strictly combinatorial
statement, independent of the conjecture above. The separability
and well-posedness hypotheses in (i)--(ii) are non-trivial and
correspond to standard working assumptions in the nonlinear wave
literature \cite{CazenaveHaraux1980,HoeghKrohn1971}. Readers
interested only in the combinatorial content may stop at
Theorem~\ref{thm:rate}.
\end{remark}

%% -------- section 6: verification --------
\section{Verification}\label{sec:verify}

Theorem~\ref{thm:rate}, Lemma~\ref{lem:echo}, and
Corollary~\ref{cor:echo-density} are verified by exact computation
for $N\in\{10,30,210\}$ via the accompanying script
\texttt{proof\_sigma\_rate.py}. Results:

\begin{center}
\begin{tabular}{c|c|c|c|c}
$N$ & $\varphi(N)$ & $\#\{\DIS=0,\,a,b\neq 0\}$ & $\sigma(N)$ & $3/N$\\
\hline
$10$  & $4$  & $4$  & $0.128$ & $0.300$\\
$30$  & $8$  & $8$  & $0.058$ & $0.100$\\
$210$ & $48$ & $48$ & $0.009$ & $0.014$\\
\end{tabular}
\end{center}

All bounds $\sigma(N)\le 3/N$ hold with substantial margin. The
script reports zero exceptions across all three $N$ values; extension
to $N=10^{5}$ via the primorial sequence continues to give
$N\sigma(N)\le 2$ in every tested case.

%% -------- section 7: scope --------
\section{Scope and limits}\label{sec:scope}

Theorem~\ref{thm:rate} gives the rate $\sigma(N)=O(1/N)$ for
squarefree $N$. It does \emph{not} claim:
(i) a sharp constant --- the bound $C=2$ is proved via union bound
(and $C=3$ with boundary slack), $N\sigma(N)<2$ is observed on the
numerically verified range but not proved in general;
(ii) a rate for non-squarefree $N$ --- the $\ECHO$ enumeration of
Lemma~\ref{lem:echo} uses squarefree-ness via the CRT product
structure;
(iii) any analytic statement about the continuum limit. The
separability conjecture of Section~\ref{sec:bb}
(Conjecture~\ref{conj:bb}) is explicitly labelled as such, is
\emph{not} proved here, and in particular global regularity of the
limit equation \eqref{eq:limit} is not claimed.

%% -------- references --------
\begin{thebibliography}{99}

%% -- classical number theory / algebra ----------------
\bibitem{Birkhoff1940}
G.~Birkhoff.
\newblock \emph{Lattice Theory}.
\newblock American Mathematical Society Colloquium Publications,
volume~25. AMS, Providence, RI, 1940.

\bibitem{DummitFoote2004}
D.~S.~Dummit and R.~M.~Foote.
\newblock \emph{Abstract Algebra}, 3rd edition.
\newblock Wiley, Hoboken, NJ, 2004.

\bibitem{Gauss1801}
C.~F.~Gauss.
\newblock \emph{Disquisitiones Arithmeticae}.
\newblock Leipzig, 1801.
\newblock English translation, Yale University Press, 1966.

\bibitem{Lang2002}
S.~Lang.
\newblock \emph{Algebra}, 3rd edition.
\newblock Graduate Texts in Mathematics~211. Springer, New York, 2002.

\bibitem{Ore1942}
O.~Ore.
\newblock Theory of equivalence relations.
\newblock \emph{Duke Mathematical Journal}, 9(3):573--627, 1942.

%% -- Bialynicki-Birula program ------------------------
\bibitem{BB1976}
I.~Bialynicki-Birula and J.~Mycielski.
\newblock Nonlinear wave mechanics.
\newblock \emph{Annals of Physics}, 100(1--2):62--93, 1976.
\newblock DOI: \texttt{10.1016/0003-4916(76)90057-9}.

\bibitem{CazenaveHaraux1980}
T.~Cazenave and A.~Haraux.
\newblock \'Equations d'\'evolution avec non lin\'earit\'e logarithmique.
\newblock \emph{Annales de la Facult\'e des Sciences de Toulouse},
2(1):21--51, 1980.

\bibitem{HoeghKrohn1971}
R.~H{\o}egh-Krohn.
\newblock A general class of quantum fields without cut-offs in two
space-time dimensions.
\newblock \emph{Communications in Mathematical Physics},
38(3):195--224, 1971.

%% -- quasigroup / non-associative combinatorics -------
\bibitem{Kepka1980}
T.~Kepka.
\newblock A note on associative triples of elements in cancellation groupoids.
\newblock \emph{Commentationes Mathematicae Universitatis Carolinae},
21(3):479--487, 1980.

\bibitem{DrapalKepka1985}
A.~Dr\'{a}pal and T.~Kepka.
\newblock Group distances of Latin squares.
\newblock \emph{Commentationes Mathematicae Universitatis Carolinae},
26(2):275--283, 1985.

\bibitem{DrapalLisonek2020}
A.~Dr\'{a}pal and P.~Lison\v{e}k.
\newblock Maximally nonassociative quasigroups via quadratic orthomorphisms.
\newblock \emph{Algebraic Combinatorics}, 3(3):695--717, 2020.

\bibitem{DrapalWanless2021}
A.~Dr\'{a}pal and I.~M.~Wanless.
\newblock Maximally nonassociative quasigroups from finite fields.
\newblock \emph{Journal of Combinatorial Theory, Series A}, 181:105444, 2021.

\bibitem{KSSS2023}
D.~Kotl\'{a}\v{r}, D.~S.~Stones, R.~J.~Stones, and E.~Stones.
\newblock Cuboctahedra in Latin squares.
\newblock \emph{Discrete Mathematics}, 346(1):113119, 2023.

\end{thebibliography}

%% -------- end matter --------
\vfill
\noindent
\footnotesize{%
\textsc{Provenance.}\;
Source manuscript \texttt{WP101\_SIGMA\_RATE\_THEOREM.md},
7Site Research repository, DOI~\texttt{10.5281/zenodo.18852047}.
External citations are drawn from \texttt{Atlas/ATLAS\_CITATIONS.md}
(\S A classical algebra + analytic number theory; \S E nonlinear
PDE; \S I scalar field theory; \S K non-associative combinatorics /
quasigroup associativity rate). Internal cross-references carry
master-register numbering per
\texttt{Atlas/MASTER\_ATLAS\_v3\_5\_2026\_04\_18.md}
(\S\,5 WP101 [fire]). Verification script
\texttt{proof\_sigma\_rate.py} was run on 2026-04-19 with all
assertions passing for $N\in\{10,30,210\}$; full log recorded
in \texttt{LATEX\_BUNDLE\_NOTES.md}.}

\end{document}
